\documentclass[11pt]{article}

\usepackage[margin=1in]{geometry}
\usepackage{hyperref}
\usepackage{amsmath}

\newcommand{\code}[1]{\texttt{#1}}

\title{\code{kv} and \code{pred}}
\author{}
\date{}

\begin{document}
\maketitle

\begin{abstract}
\code{kv} is a read-optimized key/value set with a hash signature that answers
subset queries without touching the elements. \code{pred} builds on it: a
label-selector style predicate that is compiled once and matched many times.
Both trade construction cost for match cost, on the assumption that a set or a
predicate is built once and queried repeatedly. They live at
\code{github.com/bowei/kvec/pkg/kv} and \code{github.com/bowei/kvec/pkg/pred}.
\end{abstract}

\section{\code{kv}}

\subsection{Public interface}

\paragraph{\code{KV}} An immutable set of unique key/value pairs.

\begin{description}
  \item[\code{New(itr func(yield func(k, v string) bool), cap int) *KV}]
    builds a \code{KV} from an iterator over pairs. Keys are assumed unique;
    the constructor does not deduplicate. \code{cap} is a capacity hint equal to
    the expected number of pairs.
  \item[\code{(*KV) Len() int}] the number of pairs.
  \item[\code{(*KV) Equal(other *KV) bool}] whether both sides hold exactly the
    same pairs.
  \item[\code{(*KV) Contains(other *KV) bool}] whether every pair of
    \code{other} is also in the receiver, i.e.\ whether \code{other} is a subset.
  \item[\code{(*KV) ContainsKeys(keys ...string) bool}] whether the receiver
    holds a pair for every named key, whatever the value. No keys is vacuously
    true; duplicate keys are checked more than once.
  \item[\code{(*KV) Get(key string) (string, bool)}] the value for a key, and
    whether such a pair exists.
\end{description}

\paragraph{\code{Signature}} A standalone, dynamically sized approximation of
set containment over plain strings.

\begin{description}
  \item[\code{NewSignature(bitWidth uint, keys ...string) *Signature}] rounds
    \code{bitWidth} up to a power of two, minimum $64$, and sets one bit per
    key. Duplicate keys are harmless. Panics above $2^{24}$ bits.
  \item[\code{(*Signature) BitWidth() uint}] the width actually allocated.
  \item[\code{(*Signature) Contains(other *Signature) bool}] whether every bit
    of \code{other} is set in the receiver. Panics on mismatched widths, since
    no answer would preserve the one-sided guarantee below.
  \item[\code{(*Signature) Equal(other *Signature) bool}] same width, same bits.
\end{description}

The guarantee is one-sided. For string sets $A, B$ with signatures $S_A, S_B$:
\[
  A \supseteq B \implies S_A.\mathrm{Contains}(S_B),
\]
but the converse does not hold, as distinct keys can fold onto the same bit. A
signature therefore rejects with certainty and accepts only provisionally.

\subsection{Internal optimizations}

\paragraph{Fixed-width signatures as a rejection gate.} Every \code{KV} carries
two \code{sigLen}$\,\times\,64$-bit vectors built at construction: \code{ksig}
over the key hashes, and \code{kvsig} over the key hashes mixed with the value
hashes. \code{Contains} tests \code{other}'s bits against the receiver's with
\code{\&\^{}}; any bit present in \code{other} and absent in the receiver proves
a pair the receiver cannot hold, and rejects before a single element is read.
Both vectors are checked. This is not redundant: the two fold different hashes
into the same width, so a pair missing from one can collide in the other, and at
\code{sigLen}${}=2$ both vectors sit in the cache line already fetched. The
fixed width lets the loops be unrolled and keeps the vector inline in the
struct, where the dynamic \code{Signature} would cost an allocation and an
unrollable loop.

\paragraph{Hash-ordered storage.} Pairs are sorted once at construction by
\code{kvCmp}, which orders on the 64-bit FNV-1a hash of the key and breaks ties
on the key itself so the order is total. Hash order makes the sort key an
integer comparison in the common case, and the totality is what \code{Equal},
\code{Contains}, \code{ContainsKeys} and \code{Get} all rely on for their merge
and binary searches.

\paragraph{Galloping merge in \code{Contains}.} Both sides being sorted, the
subset test is a single merge pass. A pure merge costs a comparison per element
of the larger side, which is wasteful when the receiver is much larger than
\code{other}. The merge therefore counts consecutive skipped elements, and once
the run reaches \code{gallopAfter} ($=4$) it switches to a binary search over
the remaining tail. The threshold is the point at which the walk has already
cost about what $\log_2$ comparisons would, so the search never loses by much
and wins outright on longer runs. The check sits inside the skip branch rather
than at the top of the loop, so a merge of two similar sets --- which never
builds a run that long --- does not pay for it. The pass also aborts early
whenever more elements remain in \code{other} than in the receiver.

\paragraph{Per-probe signature gate.} \code{ContainsKeys} and \code{Get} build a
one-bit signature for the probe key and test it against \code{ksig} before
searching. A missing key is usually rejected on a bit operation instead of a
binary search.

\paragraph{Allocation-free hashing.} \code{fnv1a64} is written out inline rather
than through \code{hash.Hash}, avoiding the hash-object allocation on a path
taken once per key at construction and once per probe at query time.

\paragraph{Power-of-two widths.} Both signature types require a power-of-two bit
width, so folding a hash to a bit offset is a mask (\code{hash \& bitMask})
rather than a division. \code{Signature.Contains} reslices the other side to the
receiver's length so the bounds-check eliminator can drop the checks in the
loop body.

\section{\code{pred}}

\subsection{Public interface}

\paragraph{\code{PredicateBuilder}} accumulates terms; each method returns the
receiver so calls chain.

\begingroup\sloppy
\begin{description}
  \item[\code{NewPredicateBuilder() *PredicateBuilder}] an empty builder.
  \item[\code{Key(keys ...string)}] each key must be present, whatever its value.
  \item[\code{NoKey(keys ...string)}] none of the keys may be present.
  \item[\code{Has(kv map[string]string)}] each key must be present with the
    given value. A later call overwrites an earlier one on the same key.
  \item[\code{ValueIn(key string, values []string)}] the key must be present
    with a value drawn from \code{values}. An empty \code{values} makes the
    predicate unsatisfiable.
  \item[\code{ValueNotIn(key string, values []string)}] the key must be absent,
    or present with a value outside \code{values}.
  \item[\code{Build() *Predicate}] compiles the accumulated terms. The result is
    independent of the builder, and \code{Build} may be called more than once.
\end{description}
\endgroup

\paragraph{\code{Predicate}} is the compiled, immutable conjunction of those
terms. \code{(*Predicate) Match(kv *kv.KV) bool} reports whether a \code{KV}
satisfies every term; a predicate built from an empty builder matches
everything. Being immutable, one \code{Predicate} can be matched against many
\code{KV}s from many goroutines.

\begin{verbatim}
p := NewPredicateBuilder().
    Key("app").
    NoKey("deprecated").
    Has(map[string]string{"tier": "backend"}).
    ValueIn("env", []string{"prod", "staging"}).
    Build()
\end{verbatim}

\subsection{Internal optimizations}

The builder stores terms as given and does no analysis. All of it happens once,
in \code{Build}, so that \code{Match} does the least possible work per candidate.

\paragraph{Merging repeated terms.} Repeated \code{ValueIn} terms on one key are
intersected and repeated \code{ValueNotIn} terms are unioned, so a key costs one
lookup at match time however many times it was named.

\paragraph{Subsumption.} A term that settles a key's value makes weaker terms on
that key dead, and they are dropped:
\begin{itemize}
  \item an equality decides both value tests on its key at build time, and makes
    a \code{Key} on it redundant;
  \item an in-set absorbs an exclusion on the same key by subtracting it, and
    itself proves the key present, so a \code{Key} on it is also dropped;
  \item a \code{NoKey} satisfies an exclusion on the same key vacuously, so the
    exclusion is dropped.
\end{itemize}

\paragraph{Demotion to equality.} An in-set reduced to a single value is
rewritten as an equality. That moves it behind the signature gate of
\code{kv.KV.Contains}, turning a lookup into part of a bit test.

\paragraph{Batched equalities.} All pinned pairs are compiled into one
\code{kv.KV}, so the whole group is tested by a single \code{Contains} rather
than a lookup per pair --- and rejected on two words of signature in the common
case.

\paragraph{Contradictions resolved once.} Conflicts between terms (an equality
outside its in-set, an in-set emptied by subtraction, a key required both
present and absent) are found at build time and recorded in a single
\code{unsatisfiable} flag, so \code{Match} returns on one boolean.

\paragraph{Cheapest-first evaluation.} \code{Match} orders its terms by cost:
the equality set first, since its signature gate rejects the whole group without
a lookup and it is the only test that can reject on a wrong value rather than a
missing key; then the existence tests, each usually rejecting on a signature bit
alone; then the value tests, which always pay for a lookup.

\paragraph{\code{sortedset}.} Value sets are sorted and deduplicated at build
time, which both shrinks them and enables searching. \code{contains} scans
linearly up to \code{linearScanMax} ($=8$) elements, where locality and branch
prediction beat the $\log_2 n$ comparisons of a binary search, and binary
searches above it. \code{intersect}, \code{union} and \code{subtract} are single
linear merges over the sorted representation.

\paragraph{Deterministic layout.} The per-key sets are flattened into slices
ordered by key, so a compiled \code{Predicate} does not depend on map iteration
order and matches in a stable, cache-friendly sequence.

\section{Benchmark results}

All numbers below were measured on an Apple M3 under \code{go1.25.1
darwin/arm64} with \code{go test -bench . -benchmem -count 5}; each figure is
the median of the five runs, in nanoseconds per operation. Run-to-run spread was
at most $5\%$ of the median, and under $1\%$ for most figures. Every
operation on the query path allocates nothing: only \code{Build} appears with a
non-zero allocation count.

\subsection{\code{kv.KV}}

Table~\ref{tab:kv} is a 64-pair set. The pattern the design is built around
shows up directly: a miss is more than an order of magnitude cheaper than a hit,
because the signature answers it without reaching the elements.

\begin{table}[h]
\centering
\begin{tabular}{lrl}
\hline
Operation & ns/op & Path taken \\
\hline
\code{Contains} hit (2 pairs)      &  48.6 & merge, then compare values \\
\code{Contains} miss (2 pairs)     &   1.4 & rejected on \code{kvsig} \\
\code{ContainsKeys} hit            &  32.3 & signature passes, binary search \\
\code{ContainsKeys} miss           &   6.0 & rejected on \code{ksig} \\
\code{ContainsKeys} 4 keys, all hit & 129.4 & four searches \\
\code{ContainsKeys} 5 keys, last misses & 133.1 & four searches, then reject \\
\hline
\end{tabular}
\caption{\code{KV} operations against a 64-pair set.}
\label{tab:kv}
\end{table}

\subsection{\code{Contains} across sizes}

Table~\ref{tab:pairwise} times \code{Contains} for every pair of set sizes. At
these sizes both signatures are saturated on both sides, so they never reject:
this is the case where the merge has to do the work.

\begin{table}[h]
\centering
\begin{tabular}{lrrrr}
\hline
 & \multicolumn{4}{c}{\code{other}} \\
receiver & 1k & 10k & 100k & 1M \\
\hline
1k   &      4\,667 &     1.7 &      1.8 &      1.8 \\
10k  &      6\,451 & 47\,745 &      1.7 &      1.7 \\
100k &      8\,180 & 101\,523 & 470\,658 &      1.8 \\
1M   &     16\,218 & 188\,793 & 2\,907\,257 & 4\,704\,127 \\
\hline
\end{tabular}
\caption{\code{Contains} in ns/op. Below the diagonal the call is a genuine hit
that runs the merge to completion; above it the receiver is the smaller set and
the call exits on the length check.}
\label{tab:pairwise}
\end{table}

Two things are worth reading off it. First, the upper triangle is flat at
$1.7$--$1.8$\,ns: a set that cannot possibly be a superset is rejected in
constant time, whatever the sizes involved. Second, the lower triangle is where
the gallop earns its keep. Probing 1k pairs into a 1M-pair set costs
$16.2\,\mu$s, about $16$\,ns per probe --- the cost of a binary search each.
Walking that set instead would cost roughly what the 1M/1M cell costs,
$4.7$\,ms, so on this shape the gallop is around $290\times$ ahead. The
diagonal, where the two sides are the same size and the gallop never triggers,
runs at about $4.7$\,ns per pair.

\subsection{\code{kv.Signature}}

\begin{table}[h]
\centering
\begin{tabular}{lrrr}
\hline
 & 64 bits & 1024 bits & 65536 bits \\
\hline
hit  & 1.9 & 6.4 & 283.8 \\
miss & 1.4 & 2.9 &  92.4 \\
\hline
\end{tabular}
\caption{\code{Signature.Contains} in ns/op, 64 keys in the receiver.}
\label{tab:sig}
\end{table}

Cost is linear in the number of words, since a hit has to read all of them; a
miss is cheaper because the loop returns at the first word that fails. This is
the argument for the inline fixed-width signature inside \code{KV}: at
\code{sigLen}${}=2$ the whole test is the $1.4$--$1.9$\,ns line of
Table~\ref{tab:sig}, with no allocation and no loop.

\subsection{Where the signature gate stops paying}

The gate is only worth what it rejects, and what it rejects decays as the set
fills it. A \code{KV} carries $64 \times \code{sigLen} = 128$ bits, so a set of
$n$ keys leaves a fraction of about
\[
  \left(1 - \frac{1}{128}\right)^{n}
\]
of them clear, and only an absent key landing on a clear bit is rejected without
a search. Table~\ref{tab:sat} measures both the rate and its cost. The share is
counted over 4096 absent keys; the timing rotates through 64 of them, so each
figure averages over the gate firing and not firing rather than depending on
whether one chosen key happens to collide.

\begin{table}[h]
\centering
\begin{tabular}{lrrrrrrr}
\hline
pairs in the set & 8 & 16 & 32 & 64 & 128 & 256 & 1024 \\
\hline
absent keys rejected by \code{ksig} & 92.9\% & 87.2\% & 75.3\% & 64.5\% & 31.2\% & 13.9\% & 0.0\% \\
predicted by the formula            & 93.9\% & 88.2\% & 77.8\% & 60.5\% & 36.6\% & 13.4\% & 0.0\% \\
\code{ContainsKeys} miss, ns/op     & 12.8 & 13.5 & 14.7 & 16.4 & 30.5 & 43.3 & 50.6 \\
\hline
\end{tabular}
\caption{Decay of the fixed signature as the set grows. Measured rates track the
formula to within a few points.}
\label{tab:sat}
\end{table}

The cost of a missing key rises from $12.8$ to $50.6$\,ns as the gate goes from
rejecting almost everything to rejecting nothing, at which point every probe
pays for the binary search the signature exists to avoid. This is the
quantitative form of the note in the source that \code{sigLen} suits sets of far
fewer than a thousand pairs: the gate is worth having up to a hundred or so
pairs, is half wasted by $128$, and is dead weight by $1024$. A set expected to
run large wants a larger \code{sigLen}, at two words of overhead per doubling.

Note that the single figure in Table~\ref{tab:kv} --- a $6.0$\,ns
\code{ContainsKeys} miss at 64 pairs --- is one key that happened to land on a
clear bit. Table~\ref{tab:sat} gives the honest average for that size, $16.4$\,ns.

\subsection{\code{pred}}

\code{Match} checks its terms cheapest-first, so a single ``miss'' figure would
say nothing: it would report whichever term the input happened to break.
Table~\ref{tab:match} instead breaks a different term in each column, in the
order \code{Match} tries them.

\begin{table}[h]
\centering
\begin{tabular}{llrrr}
\hline
\multicolumn{2}{l}{rejected on} & \multicolumn{3}{c}{\code{KV} size} \\
term & broken by & 8 & 16 & 64 \\
\hline
\code{eq}         & key absent        &   2.4 &   2.4 &   2.4 \\
\code{eq}         & wrong value       &   2.4 &   2.4 &   2.4 \\
\code{keyExists}  & required key gone  &  45.2 &  66.7 &  77.9 \\
\code{keyNExists} & forbidden key present &  93.8 & 104.1 & 122.5 \\
\code{in}         & value outside the set &  91.2 & 103.9 & 120.9 \\
\code{notIn}      & value inside the set & 116.4 & 132.7 & 155.6 \\
\hline
\multicolumn{2}{l}{accepted (all terms pass)} & 117.1 & 133.3 & 157.2 \\
\multicolumn{2}{l}{unsatisfiable predicate}   & \multicolumn{3}{c}{1.9} \\
\hline
\end{tabular}
\caption{\code{Match} in ns/op, for a five-term predicate naming one term of
every kind.}
\label{tab:match}
\end{table}

The spread is the result. A candidate rejected on the equality set costs
$2.4$\,ns whatever its size, against $157.2$\,ns for one that passes every term
--- a factor of $65$ between the cheapest and dearest outcome of the same call.
Ordering the terms is worth that factor only because the cheap tests are also
the ones most likely to fire: a selector applied across a population rejects
most of it, and rejects most of that on a pinned pair.

Three details are worth reading carefully.

\begin{itemize}
  \item A wrong value costs exactly what an absent key costs, $2.4$\,ns. The
    \code{eq} test is not a lookup that happens to fail early --- \code{kvsig}
    folds the value into the bit, so a pair with the right key and the wrong
    value misses the signature just as a missing key does.
  \item The rows are not a strict ladder, and cannot be: breaking a term also
    changes what the earlier terms do. Rejecting on \code{keyNExists} requires
    the forbidden key to be \emph{present}, which turns the lookup that rejects
    into a hit rather than a signature miss, and that is why it costs about what
    the later \code{in} row costs. Each row is the true cost of that outcome,
    but the rows are not a decomposition of one call.
  \item Everything on this path is allocation-free at every size.
\end{itemize}

\code{Build} costs $953.9$\,ns and $1584$\,B in $30$ allocations for this
predicate. That is the bet of the package as a number: against the $157.2$\,ns
accept, it pays for itself after about six matches, and against the $2.4$\,ns
reject, after roughly four hundred. Either way a predicate matched across a
population of any size has long since repaid it.

The same five terms written so that two of them are redundant --- a \code{Key}
and a \code{ValueNotIn} on a key that an equality already settles --- build in
$837.0$\,ns and $25$ allocations, and compile to three terms rather than five.
The redundant form is both cheaper to build and cheaper to match, since the work
it names is discarded rather than performed once per candidate.

\subsection{Why the equality set is one \code{Contains}}

\code{BenchmarkEqStrategy} compares the two ways of testing three pinned pairs:
the single \code{Contains} that \code{Match} uses, against the lookup-per-pair
it replaced.

\begin{table}[h]
\centering
\begin{tabular}{llrrr}
\hline
 & \code{KV} size & 8 & 16 & 64 \\
\hline
hit  & \code{Contains}  & 41.8 & 47.2 & 55.3 \\
     & per-pair lookup  & 99.0 & 109.0 & 130.6 \\
\hline
miss & \code{Contains}  &  1.9 &  1.9 &  1.9 \\
     & per-pair lookup  & 83.9 & 92.6 & 109.0 \\
\hline
\end{tabular}
\caption{Testing three pinned pairs, in ns/op.}
\label{tab:eq}
\end{table}

\code{Contains} is $2.3$--$2.4\times$ faster on a hit, where it makes one merge
pass instead of three independent searches. On a miss it is $45$--$58\times$
faster,
and --- more to the point --- flat at $1.9$\,ns across every size, where the
lookup form grows with the candidate. That flat line is the $2.4$\,ns
\code{eq} row of Table~\ref{tab:match} seen on its own, and it is the reason
\code{Build} folds every equality into a single \code{kv.KV} and \code{Match}
tests it first.

\end{document}
